\section{Related Work}
\label{sec:related}

\subsection{Research Artifacts and Scholarly Records}
Modern open-science and digital-library research increasingly treats datasets, software, URLs, code repositories, and packaged research objects as core elements of the broader scholarly record~\cite{wilkinson2016fair,datacitation2014joint,smith2016softwarecitation,lamprecht2020fairsoftware,katz2021freshfairsoftware,barker2022fair4rs,soilandreyes2022rocrate}.
Frameworks such as the FAIR data principles, software citation standards, FAIR research software, and RO-Crate protocols emphasize that computational artifacts require persistent identifiers, explicit authorship, licenses, dependencies, and contextual metadata instead of simple access hyperlinks~\cite{wilkinson2016fair,datacitation2014joint,smith2016softwarecitation,lamprecht2020fairsoftware,katz2021freshfairsoftware,barker2022fair4rs,soilandreyes2022rocrate}.
Analysis of scholarly URLs similarly reveals that a visible hyperlink falls short of a functional curation record because its surrounding textual context must be interpreted before the link can support metadata generation, discovery, or long-term preservation choices~\cite{salsabil2025context}.
Existing scholarship establishes the bibliographic value of detailed artifact descriptions while primarily focusing on established artifact classes such as standalone datasets, software packages, URLs, or structured research objects.
Distributed record fragmentation becomes more pronounced for compact or derived open \llm artifacts because their operational identity, scholarly relationships, upstream dependencies, and release components may be scattered across papers, hub repositories, model cards, code repositories, and developer statements~\cite{mitchell2019modelcards,pepe2024hfdocs,kadasi2025modelhubs,jiang2023ptmreuse,taraghi2024modelreuse}.
Building on these foundations, our study asks when fragmented public traces mature into sufficient evidence for institutional digital-library curation.

\subsection{Reuse, Reproducibility, and Release Packages}
Existing literature on repository reuse and computational reproducibility draws a sharp distinction between baseline availability and functional reuse~\cite{shiri2023reuse,stodden2014implementing,akella2023effort,dodge2019show,pineau2021reproducibility}.
Investigations into repository interfaces demonstrate that evaluating reuse requires tracking how digital objects are shared, transformed, and recontextualized~\cite{shiri2023reuse}.
Methodological studies and reproducibility reporting programs similarly show that publishing raw source code does not reflect the full effort required to inspect, rerun, or validate computational findings; auxiliary artifacts, documentation, environments, and procedural details also matter~\cite{stodden2014implementing,akella2023effort,dodge2019show,pineau2021reproducibility}.
This broader view motivates evaluating release evidence as a multi-layered package rather than classifying availability through a simple open-or-closed variable.
Emerging open \llm variants introduce distribution patterns that are not well captured by generic software availability fields: a single record may expose training code, fine-tuning recipes, checkpoint weights, LoRA adapters, quantized files, evaluation assets, or only an unfulfilled statement of intent~\cite{jiang2023ptmreuse,taraghi2024modelreuse,white2024mof}.
Grounded in this multi-tiered perspective, our release-package hierarchy categorizes which visible assets a publication record exposes for inspection and curation without assuming that those assets are executable or complete.

\subsection{Provenance and Derived-Model Relations}
Provenance frameworks have long been foundational to digital-library architectures, scientific workflows, and computational reproducibility~\cite{moreau2013prov,davidson2008provenance,herschel2017survey}.
Standardized representations such as the W3C PROV model formalize lineage by mapping entities, activities, and agents to explain how research artifacts are derived and produced~\cite{moreau2013prov}.
Digital-library systems further illustrate how provenance tracking can be integrated into document-analysis workflows so that generated outputs, processing steps, and analytical decisions remain traceable~\cite{rauch2025ontextract}.
Open \llm deployments introduce a related but more relation-specific challenge.
A derived or compressed model variant may depend on an upstream artifact as a base checkpoint, teacher for distillation, adapter target, merge source, quantization source, evaluator, or baseline~\cite{jiang2023ptmreuse,taraghi2024modelreuse,white2024mof}.
Available public records frequently mention broad model-family names or contextual references that aid basic discovery without encoding an explicit, machine-readable upstream relationship.
Addressing this limitation requires treating provenance as a tiered spectrum of public evidence, where broad family-level signals support catalog organization while rigorous curation demands canonical upstream declarations and explicit relationship types.

\subsection{Model Documentation, Model Hubs, and Open AI Infrastructure}
Documentation paradigms such as model cards and structured datasheets provide influential templates for recording the properties of machine-learning models and datasets~\cite{mitchell2019modelcards,crisan2022interactive,shen2022modelcardtoolkit,gebru2021datasheets}.
Frameworks such as Data Cards and the Croissant format extend this documentation agenda by encoding data assets and ML-ready resources as structured metadata objects~\cite{pushkarna2022datacards,akhtar2024croissant}.
Quantitative analyses of active model repositories reveal that documentation practices vary substantially across communities, particularly in how authors report licensing, training datasets, model properties, and reuse guidance~\cite{liang2024modelcards,pepe2024hfdocs,kadasi2025modelhubs}.
Independent tracking of \hf repository activity further demonstrates how pretrained models undergo reuse, adaptation, and open-source distribution throughout AI developer communities~\cite{jiang2023ptmreuse,taraghi2024modelreuse,osborne2024aicommunity}.
Collaborative transparency frameworks similarly emphasize that openness depends on multiple technical components rather than model weights alone~\cite{bommasani2023fmi,white2024mof}.
Prior landscape evaluations map how infrastructure platforms distribute and govern AI artifacts, yet they typically isolate a single documentation layer such as a model card, repository field, reuse network, or openness checklist.
Our research shifts the unit of analysis toward the cross-system scholarly record.
Connecting repository metadata with scholarly publications allows us to ask what existing public records reveal about artifact identity, paper-to-repository linkage, upstream-evidence claims, inspectable release assets, and integrated curation readiness.
